\section{Resumo por Estrofes}

Nesta seção, apresenta-se um resumo das cinquenta primeiras estrofes de
\citetitle{camoes_lusiadas}\cite{camoes_lusiadas}, destacando os principais
acontecimentos, personagens e elementos narrativos de cada trecho. O objetivo é
oferecer uma visão geral do desenvolvimento da narrativa antes de sua análise,
permitindo ao leitor acompanhar com maior clareza os episódios e temas que
serão posteriormente aprofundados.

\medskip

\textbf{1--3:}\\
Anúncio do que será cantado: os feitos portugueses.

\medskip

\textbf{4--5:}\\
Pede inspiração às Tágides (ninfas do rio Tejo) para escrever seu poema.

\medskip

\textbf{6--8:}\\
Dedicatória a Dom Sebastião e apresentação dele como esperança de Portugal.

\medskip

\textbf{9--18:}\\
Continua a dedicatória aos heróis portugueses e à grandeza de Portugal,
desejando que o rei continue essa tradição.

\medskip

\textbf{19:}\\
As naus portuguesas já estão navegando; começa a narrativa da viagem.

\medskip

\textbf{20--22:}\\
Concílio dos deuses: agora, no plano mitológico, os deuses se reúnem para
decidir o destino dos portugueses.

\medskip

\textbf{23--30:}
\begin{itemize}
    \item Júpiter apoia os portugueses;
    \item o Fado (destino) determina que eles alcancem o Oriente;
    \item Júpiter decide apoiá-los e permitir seu acolhimento na África;
    \item Baco se opõe aos portugueses com medo de perder sua fama no Oriente.
\end{itemize}

\medskip

\textbf{31--36:}
\begin{itemize}
    \item Baco é contra os portugueses;
    \item Vênus é a favor, admira suas características e deseja sua própria
      glória;
    \item Marte apoia Vênus e reconhece o valor guerreiro lusitano.
\end{itemize}

\medskip

\textbf{36--41:}
\begin{itemize}
    \item Marte intervém e critica a inveja de Baco;
    \item Júpiter mantém sua decisão;
    \item Mercúrio será enviado para ajudar os portugueses;
    \item os deuses encerram o concílio.
\end{itemize}

\medskip

\textbf{42--43:}\\
A navegação é tranquila.

\medskip

\textbf{44:}\\
Vasco da Gama é introduzido: corajoso e determinado, não vê motivo para
interromper a expedição.

\medskip

\textbf{45--48:}\\
Encontro com os habitantes locais e descrição de seus costumes, aparência,
vestimentas e embarcações.

\medskip

\textbf{49--50:}\\
Primeiro contato e diálogo entre os povos, marcado por curiosidade mútua.
